\subsection{Reading mathematical notation}
\label{math:language}

\subsubsection{Indices, sums, and products}

An index labels members of a collection.
In $x_i$, the index $i$ might label a cell, measurement, spatial site, or component of a vector.
The summation notation
\[
    \sum_{i=1}^{N}x_i=x_1+x_2+\cdots+x_N
\]
means to add the terms obtained by letting $i$ run from $1$ to $N$.
Similarly,
\[
    \prod_{i=1}^{N}x_i=x_1x_2\cdots x_N.
\]

Indices are dummy labels inside a complete sum.
For example, $\sum_{i=1}^{N}x_i=\sum_{j=1}^{N}x_j$.
The symbol used for the index changes, but the value does not.

\paragraph{Worked example: a sample mean}
For measurements $x_1=2$, $x_2=5$, and $x_3=8$, the sample mean is
\[
    \bar{x}=\frac{1}{3}\sum_{i=1}^{3}x_i=\frac{2+5+8}{3}=5.
\]

\subsubsection{Powers, exponentials, and logarithms}

The exponential function $e^x$ is special because its derivative equals itself.
It converts addition in the exponent into multiplication:
\[
    e^{a+b}=e^ae^b.
\]
The natural logarithm $\ln x$ is its inverse:
\[
    \ln(e^x)=x, \qquad e^{\ln x}=x \quad (x>0).
\]
Useful logarithm rules are
\[
    \ln(xy)=\ln x+\ln y, \qquad \ln\left(\frac{x}{y}\right)=\ln x-\ln y, \qquad \ln(x^a)=a\ln x.
\]

Exponential growth has the form $M(t)=M_0e^{kt}$ with $k>0$.
Its doubling time $t_2$ satisfies $M(t_2)=2M_0$, so
\[
    t_2=\frac{\ln2}{k}.
\]
Exponential decay has the form $A(t)=A_0e^{-t/\tau}$.
The time $\tau$ is the time at which the amplitude has fallen to $1/e$ of its initial value.



